\documentclass[11pt]{article}
\usepackage[a4paper,margin=25mm]{geometry}
\usepackage[no-math]{fontspec}
\setmainfont{Latin Modern Roman}
\usepackage{amsmath,amssymb,amsthm,xfrac,xspace,hyperref}
\newcommand{\ie}{i.e.,\xspace}
\newcommand{\eg}{e.g.,\xspace}
\newcommand{\etc}{etc.\xspace}
\newcommand{\cf}{cf.\xspace}
\newcommand{\Subsection}{\subsection}
\newcommand{\Subsubsection}{\subsubsection}
\newcommand{\skpt}{\smallskip}
\input{author-preamble.tex}
\emergencystretch=3em
\begin{document}
\begin{center}{\Large Stable item 1920}\end{center}
\paragraph{Source and attribution.}
Eduardo Oregón-Reyes, \emph{A new inequality about matrix products and a Berger-Wang formula}, Journal de l’École polytechnique — Mathématiques 7 (2020), publisher version, DOI 10.5802/jep.114. \href{https://jep.centre-mersenne.org/articles/10.5802/jep.114/}{Publisher article}. Authors' text: \href{https://creativecommons.org/licenses/by/4.0/}{CC BY 4.0}.
\paragraph{Editorial source boundary.}
Prove only the fixed-length \(n=N(d)\) branch of the authors' Theorem 1.2. For a local field \(k\), dimension \(d\ge1\), and a submultiplicative norm on \(M_d(k)\), the exact inequality, constant bounds, and zero-matrix convention are reproduced in Theorem 1.2 below; the selected case sets the product length equal to \(N\). The authors' exterior-power/rank-drop argument, effective Nullstellensatz construction, coefficient counting and local-field estimates are retained in full as the supporting core. The article's \(n>N\) extension, Berger--Wang consequences and low-dimensional examples are preserved as context and are not selected conclusions.
\paragraph{Difficulty (editorial): H2.}
Exterior powers force rank reduction for blocks whose every contiguous subproduct is nilpotent. An effective multihomogeneous Nullstellensatz identity converts this zero-locus rigidity into a quantitative spectral-radius inequality respecting product order over arbitrary local fields.
\paragraph{Transcription note (editorial).}
Author mathematical text and ordering are unchanged. Only publisher front matter is replaced by this independent wrapper. Original TeX body and macros remain separate editable inputs. Bibliography entries are recovered from the same cached publisher article HTML. No mathematical repair or assistant-generated proof is supplied.
\clearpage
\input{author-body.tex}
\begin{thebibliography}{99}
\bibitem{avibochi} A. Avila \& J. Bochi - “A formula with some applications to the theory of Lyapunov exponents” , Israel J. Math. 131 (2002), p. 125-137
\bibitem{bw} M. A. Berger \& Y. Wang - “Bounded semigroups of matrices” , Linear Algebra Appl. 166 (1992), p. 21-27
\bibitem{bochilya} J. Bochi - “Genericity of zero Lyapunov exponents” , Ergodic Theory Dynam. Systems 22 (2002) no. 6, p. 1667-1696
\bibitem{bociineq} J. Bochi - “Inequalities for numerical invariants of sets of matrices” , Linear Algebra Appl. 368 (2003), p. 71-81
\bibitem{bro} E. Breuillard - “A height gap theorem for finite subsets of \(\mathrm{GL}_d(\overline{\mathbb Q})\) and nonamenable subgroups” , Ann. of Math. (2) 174 (2011) no. 2, p. 1057-1110
\bibitem{bri} M. R. Bridson \& A. Haefliger - Metric spaces of non-positive curvature , Grundlehren Math. Wiss. , vol. 319 , Springer-Verlag, Berlin, 1999
\bibitem{papa} M. Coornaert, T. Delzant \& A. Papadopoulos - Géométrie et théorie des groupes. Les groupes hyperboliques de Gromov , Lect. Notes in Math. , vol. 1441 , Springer-Verlag, Berlin, 1990
\bibitem{geodyn} T. Das, D. Simmons \& M. Urbański - Geometry and dynamics in Gromov hyperbolic metric spaces. With an emphasis on non-proper settings , Math. Surveys and Monographs , vol. 218 , American Mathematical Society, Providence, RI, 2017
\bibitem{elsner} L. Elsner - “The generalized spectral-radius theorem: an analytic-geometric proof” , Linear Algebra Appl. 220 (1995), p. 151-159, Proceedings of the Workshop “Nonnegative Matrices, Applications and Generalizations” and the Eighth Haifa Matrix Theory Conference (Haifa, 1993)
\bibitem{fulton} W. Fulton - Algebraic curves. An introduction to algebraic geometry , Advanced Book Classics , Addison-Wesley Publishing Company, Redwood City, CA, 1989
\bibitem{fisicos} I. Goldhirsch, P.-L. Sulem \& S. A. Orszag - “Stability and Lyapunov stability of dynamical systems: a differential approach and a numerical method” , Phys. D 27 (1987) no. 3, p. 311-337
\bibitem{goukar} S. Gouëzel \& A. Karlsson - “Subadditive and multiplicative ergodic theorems” , J. Eur. Math. Soc. (JEMS) (to appear)
\bibitem{gur} L. Gurvits - “Stability of discrete linear inclusion” , Linear Algebra Appl. 231 (1995), p. 47-85
\bibitem{jun} R. Jungers - The joint spectral radius. Theory and applications , Lect. Notes in Control and Information Sci. , vol. 385 , Springer-Verlag, Berlin, 2009
\bibitem{karmar} A. Karlsson \& G. A. Margulis - “A multiplicative ergodic theorem and nonpositively curved spaces” , Comm. Math. Phys. 208 (1999) no. 1, p. 107-123
\bibitem{kozyakin} V. Kozyakin - “The Berger-Wang formula for the Markovian joint spectral radius” , Linear Algebra Appl. 448 (2014), p. 315-328
\bibitem{klpasom} T. Krick, L. M. Pardo \& M. Sombra - “Sharp estimates for the arithmetic Nullstellensatz” , Duke Math. J. 109 (2001) no. 3, p. 521-598
\bibitem{langalg} S. Lang - Algebra , Graduate Texts in Math. , vol. 211 , Springer-Verlag, New York, 2002
\bibitem{lorentz} F. Lorenz - Algebra. Vol. II: Fields with structure, algebras and advanced topics , Universitext , Springer, New York, 2008
\bibitem{mcnaugh} R. McNaughton \& Y. Zalcstein - “The Burnside problem for semigroups” , J. Algebra 34 (1975), p. 292-299
\bibitem{morrisbw} I. D. Morris - “The generalised Berger-Wang formula and the spectral radius of linear cocycles” , J. Funct. Anal. 262 (2012) no. 3, p. 811-824
\bibitem{mathermorris} I. D. Morris - “Mather sets for sequences of matrices and applications to the study of joint spectral radii” , Proc. London Math. Soc. (3) 107 (2013) no. 1, p. 121-150
\bibitem{morrispr} I. D. Morris - “An inequality for the matrix pressure function and applications” , Adv. Math. 302 (2016), p. 280-308
\bibitem{yo} E. Oregón-Reyes - “Properties of sets of isometries of Gromov hyperbolic spaces” , Groups Geom. Dyn. 12 (2018) no. 3, p. 889-910
\bibitem{raro} H. Radjavi \& P. Rosenthal - Simultaneous triangularization , Universitext , Springer-Verlag, New York, 2000
\bibitem{rota} G.-C. Rota \& G. Strang - “A note on the joint spectral radius” , Indag. Math. 22 (1960), p. 379-381
\bibitem{schr} R. Schneider - Convex bodies: the Brunn-Minkowski theory , Encyclopedia of Math. and its Appl. , vol. 151 , Cambridge University Press, Cambridge, 2014
\bibitem{som} M. Sombra - “A sparse effective Nullstellensatz” , Adv. in Appl. Math. 22 (1999) no. 2, p. 271-295
\bibitem{tsblo} J. N. Tsitsiklis \& V. D. Blondel - “The Lyapunov exponent and joint spectral radius of pairs of matrices are hard—when not impossible—to compute and to approximate” , Math. Control Signals Systems 10 (1997) no. 1, p. 31-40, Correction: Ibid. , no. 4, p. 381
\end{thebibliography}
\end{document}
